
%% NOTE: the chapter, part and section heads are set by `titlesec' further
%% down, so no chapter-style package (fncychap and friends) is loaded here:
%% the two redefine \@makechapterhead and would fight over it.

\usepackage{emoji}
\usepackage{listingsutf8}

\usepackage{textcomp}
\usepackage[T1]{fontenc}
\usepackage{tgbonum}

%% Page geometry.
%%
%% The text block used to be 18.4cm wide on a 21.5cm page, i.e. about 110
%% characters per line -- roughly half again the 65-75 a reader's eye tracks
%% comfortably, and the reason the body text read like a web page rather than a
%% book. It is now 15cm, about 85 characters at 11pt Termes, with the classic
%% 2:3 inner-to-outer margin ratio of a two-sided book.
\usepackage[paperwidth=21.5cm, paperheight=29.1cm,
            textwidth=15cm, textheight=23cm,
            marginratio={2:3,2:3},
            headsep=16pt, footskip=30pt,
            twoside]{geometry}

\usepackage[dvipsnames]{xcolor}

\usepackage{multirow, booktabs}
\usepackage{multicol}

\usepackage{pmboxdraw}
\usepackage{framed,color,verbatim}

\usepackage[flushleft]{threeparttable}

\usepackage{tikzpagenodes}
\usepackage{everypage}
\usetikzlibrary{calc}

\usepackage{graphicx, float}
\usepackage{adjustbox}

\definecolor{shadecolor}{rgb}{.9, .9, .9}
\definecolor{steelblue}{rgb}{0.27, 0.51, 0.71}
\definecolor{uclagold}{rgb}{1.0, 0.7, 0.0}
\definecolor{silver}{rgb}{0.75, 0.75, 0.75}

\usepackage{fancyvrb}

\fvset{frame=single,framesep=1mm,fontfamily=courier,fontsize=\normalsize,numbers=left,framerule=.1mm,numbersep=1mm,commandchars=\\\{\}}

\usepackage[most]{tcolorbox}
\tcbset{
colback=gray!40,
colframe=white,
width=\dimexpr\textwidth\relax,
enlarge left by=0mm,
boxsep=5pt,
arc=0pt,outer arc=0pt,
}


\definecolor{applegreen}{rgb}{0.55, 0.71, 0.0}
\definecolor{antiquebrass}{rgb}{0.8, 0.58, 0.46}
\definecolor{PigBlue}{RGB}{14, 33, 160}
\definecolor{purple2}{RGB}{77, 60, 119}
\definecolor{background}{rgb}{1.0, 0.96, 0.89}


\makeatletter

\usepackage{textcomp}
%% NOTE: do *not* load `times' here: it selects the legacy Type1 `ptm' family,
%% whose 8-bit encoding has no glyph for U+2019 ('), U+2014 (---) or U+2011
%% (non-breaking hyphen), all of which occur in the prose. Under lualatex the
%% Unicode-complete Times clone (TeX Gyre Termes) is selected via fontspec
%% below instead.

%% Home-directory tilde in shell prompts. Must be the ASCII tilde: the mono
%% font used for listings has no glyph for U+02F7 (\texttildelow).
\newcommand{\mtilde}[0]{\raisebox{0.5ex}{\textasciitilde}}
\newcommand{\mtildeascii}[0]{\raisebox{0.5ex}{\textasciitilde}}

\usepackage{fontspec}

%% Times-compatible metrics, but with full Unicode coverage.
\setmainfont{TeX Gyre Termes}
\setsansfont{TeX Gyre Heros}

\newfontfamily\japan{IPAMincho}
\newfontfamily\korean{Noto Sans Mono CJK KR}
\newfontfamily\emojiF{Twitter Color Emoji}[Renderer=Harfbuzz] % emoji font

\usepackage{listings}

%% ---------------------------------------------------------------------------
%% Listing badges.
%%
%% The language of a listing is given by the background color of the listing.
%% That distinction disappears in a black and white print, and is hard to read
%% for a color blind reader, so every listing also carries a small badge in its
%% top right corner, made of a glyph and the name of the language.
%%
%% NOTE: `@' is already a letter here (\makeatletter above), so this block must
%% not close it with \makeatother.

%% Badge of the current listing, set by the `badge' key of its style.
\def\ymir@badge{}
\lst@Key{badge}{}{\def\ymir@badge{#1}}

%% A sheet of paper bearing a letter, in the style of the file icons used by
%% the figures of chapter 1. #1 is the fill color, #2 the letter.
\newcommand{\ymir@sheeticon}[2][white]{%
  \begin{tikzpicture}[baseline=-0.3ex]
    \filldraw[fill=#1, draw=black!55, line width=0.3pt]
    (0,0) -- (0,1.9ex) -- (1.05ex,1.9ex) -- (1.5ex,1.45ex) -- (1.5ex,0) -- cycle;
    \draw[draw=black!55, line width=0.3pt]
    (1.05ex,1.9ex) -- (1.05ex,1.45ex) -- (1.5ex,1.45ex);
    \node[inner sep=0pt, text=black!75] at (0.72ex,0.72ex)
    {\fontsize{5}{5}\selectfont\sffamily\bfseries #2};
  \end{tikzpicture}%
}

%% A terminal window showing a prompt.
\newcommand{\ymir@shellicon}{%
  \begin{tikzpicture}[baseline=-0.3ex]
    \filldraw[fill=black!65, draw=black!65, line width=0.3pt]
    (0,0) rectangle (1.8ex,1.9ex);
    \node[inner sep=0pt, text=white] at (0.9ex,0.9ex)
    {\fontsize{5}{5}\selectfont\ttfamily \$\_};
  \end{tikzpicture}%
}

%% A warning sign, for a listing whose code is deliberately wrong and does not
%% compile.
\newcommand{\ymir@erroricon}{%
  \begin{tikzpicture}[baseline=-0.3ex]
    \filldraw[fill=highlightcolor, draw=black!55, line width=0.3pt,
    line join=round]
    (0,0) -- (2.1ex,0) -- (1.05ex,1.9ex) -- cycle;
    \node[inner sep=0pt, text=black!75] at (1.05ex,0.55ex)
    {\fontsize{5}{5}\selectfont\sffamily\bfseries !};
  \end{tikzpicture}%
}

\newcommand{\ymir@badgelabel}[2]{%
  #1\,\textsf{\fontsize{7}{7}\selectfont\textcolor{black!60}{#2}}%
}

\newcommand{\ymirbadgeymir}{\ymir@badgelabel{\ymir@sheeticon[background!40]{Y}}{Ymir}}
\newcommand{\ymirbadgemyil}{\ymir@badgelabel{\ymir@sheeticon[teal!15]{M}}{M-YIL}}
\newcommand{\ymirbadgelyil}{\ymir@badgelabel{\ymir@sheeticon[olive!15]{L}}{L-YIL}}
\newcommand{\ymirbadgeshell}{\ymir@badgelabel{\ymir@shellicon}{Terminal}}
\newcommand{\ymirbadgeerror}{\ymir@badgelabel{\ymir@erroricon}{Invalid Ymir}}
\newcommand{\ymirbadgec}{\ymir@badgelabel{\ymir@sheeticon[purple!12]{C}}{C}}

%% Typeset the badge of the current listing, flush with the right edge of the
%% band of color, on the line above it. It occupies no vertical space, so the
%% listing itself is not moved. Inline listings (\token and friends) are left
%% alone: the badge only makes sense for a display listing.
\newcommand{\ymir@typesetbadge}{%
  \lst@ifdisplaystyle
    \ifx\ymir@badge\@empty\else
      \noindent\hb@xt@\dimexpr\linewidth+\lst@xleftmargin\relax{\hfil
        \raisebox{0.2ex}[\z@][\z@]{\ymir@badge}}%
      \par\nointerlineskip
    \fi
  \fi
}

\lst@AddToHook{Init}{\ymir@typesetbadge}
%% ---------------------------------------------------------------------------

\lstdefinestyle{coloredverbatim}{
badge=\ymirbadgeymir,
backgroundcolor=\color{background!40}, % Background color
numbers=left,                    % Position of line numbers
numberstyle=\tiny\color{purple!80!white},   % Style of line numbers
numbersep=1em,                   % Space between code and line numbers
xleftmargin=0.05\linewidth,                % Left margin
xrightmargin=-2pt,
belowskip=10pt,
aboveskip=10pt,
frame=l,                    % Frame around the code
rulecolor=\color{black!40},
framerule=0.1pt,
breaklines=true,                 % Allow line breaks
showstringspaces=false,          % Don't show spaces in strings
linewidth=0.95\linewidth,            % Fixed width size
basicstyle=\ttfamily\footnotesize, % Font style and size
captionpos=t,                     % Caption position
literate={π}{{$\pi$}}1
{Ω}{{\ensuremath{\Omega}}}1
{🙂}{\emojiF \emoji{slightly-smiling-face}}1
%% `~' is the Ymir concatenation operator: render it as an ASCII tilde.
%% \texttildelow would map to U+02F7, which the mono font does not provide.
{~}{\mtildeascii}1,
language=C,
keywordstyle=[1]\color{purple},       % keyword style
morekeywords=[1]{let, is, fn, unsafe, class, record, spawn, dg, trait, in, throw, throws, match, over, pub, prv, prot, self, cte, of, in, future, impl, macro, assert, atomic, \_\_version, \_\_pragma, mod, use, ref, lazy, def, with, exit, success, failure, catch, \_\_test, loop, panic, entity, real}, % if you want to add more keywords to the set
keywordstyle=[2]\color{PigBlue!80},
keywords=[2]{i32, f32, i8, i16, i64, isize, u8, u16, u32, u64, usize, fsize, f32, f64, f80, fsize, c8, c16, c32, Ok, Err, false, true, typeof, typeof, bool},
keywordstyle=[3]\color{purple2!80!white},
keywords=[3]{expand, alias, copy, dcopy, mut, dmut, none, null, @overlaid, move, cast},
stringstyle = \color{black!30!applegreen}
}

%% C source shown for interoperability: same presentation, but the badge says
%% the listing is not Ymir -- which is also what keeps it out of the listing
%% checker (see tools/check_listings.py).
\lstdefinestyle{cVerb}{
style=coloredverbatim,
badge=\ymirbadgec
}

%% Ymir code that is deliberately wrong: same presentation as an ordinary Ymir
%% listing, but the badge announces that the code does not compile.
\lstdefinestyle{coloredverbatimError}{
style=coloredverbatim,
badge=\ymirbadgeerror
}

%% NOTE: there used to be a third Ymir style here, `coloredverbatimCorrect',
%% used by two listings that carry a green correction highlight. It was a
%% verbatim copy of `coloredverbatim' that had since drifted: it still
%% highlighted `struct' and `@union', the spellings that became `record' and
%% `@overlaid', and it had missed twenty keywords added to the main style
%% afterwards. Since the correction highlight comes from `escapechar' at the
%% call site and not from the style, the two listings simply use
%% `coloredverbatim' and the fork is gone.

\lstdefinestyle{myilVerb}{
badge=\ymirbadgemyil,
backgroundcolor=\color{teal!10}, % Background color
numbers=left,                    % Position of line numbers
numberstyle=\tiny\color{purple!80!white},   % Style of line numbers
numbersep=1em,                   % Space between code and line numbers
xleftmargin=0.05\linewidth,                % Left margin
xrightmargin=-2pt,
aboveskip=10pt,
belowskip=10pt,
frame=l,                    % Frame around the code
rulecolor=\color{black!40},
framerule=0.1pt,
breaklines=true,                 % Allow line breaks
showstringspaces=false,          % Don't show spaces in strings
linewidth=0.95\linewidth,            % Fixed width size
basicstyle=\ttfamily\footnotesize, % Font style and size
captionpos=t,                     % Caption position
literate={π}{{$\pi$}}1
{Ω}{{\ensuremath{\Omega}}}1      
{~}{\mtildeascii}1,
language=C,
keywordstyle=[1]\color{purple},       % keyword style
morekeywords=[1]{let, is, fn, unsafe, class, struct, in, throw, throws, match, over, pub, prv, prot, self, cte, mod, use, exit, success, failure, catch, \_\_test, def, lazy, ref, loop, panic}, % if you want to add more keywords to the set
keywordstyle=[2]\color{PigBlue!100},
keywords=[2]{i32, f32, i8, i16, i64, isize, u8, u16, u32, u64, usize, fsize, f32, f64, c8, c16, c32, Ok, Err, false, true},
keywordstyle=[3]\color{purple!70!white},
keywords=[3]{expand, alias, copy, dcopy, mut, dmut, none, null, @union},
stringstyle = \color{black!30!applegreen}
}

\lstdefinestyle{lyilVerb}{
badge=\ymirbadgelyil,
backgroundcolor=\color{olive!10}, % Background color
numbers=left,                    % Position of line numbers
numberstyle=\tiny\color{purple!80!white},   % Style of line numbers
numbersep=1em,                   % Space between code and line numbers
xleftmargin=0.05\linewidth,                % Left margin
xrightmargin=-2pt,
aboveskip=10pt,
belowskip=10pt,
frame=l,                    % Frame around the code
rulecolor=\color{black!40},
framerule=0.1pt,
breaklines=true,                 % Allow line breaks
showstringspaces=false,          % Don't show spaces in strings
linewidth=0.95\linewidth,            % Fixed width size
basicstyle=\ttfamily\footnotesize, % Font style and size
captionpos=t,                     % Caption position
literate={π}{{$\pi$}}1
{Ω}{{\ensuremath{\Omega}}}1      
{~}{\mtildeascii}1,
language=C,
keywordstyle=[1]\color{purple},       % keyword style
morekeywords=[1]{let, is, fn, unsafe, class, struct, in, throw, throws, match, over, pub, prv, prot, self, cte, mod, use, exit, success, failure, catch, \_\_test, def, lazy, ref, loop, panic}, % if you want to add more keywords to the set
keywordstyle=[2]\color{PigBlue!100},
keywords=[2]{i32, f32, i8, i16, i64, isize, u8, u16, u32, u64, usize, fsize, f32, f64, c8, c16, c32, Ok, Err, false, true},
keywordstyle=[3]\color{purple!70!white},
keywords=[3]{expand, alias, copy, dcopy, mut, dmut, none, null, @union},
stringstyle = \color{black!30!applegreen}
}

\lstdefinestyle{bashVerb}{
badge=\ymirbadgeshell,
numbers=left,                    % Position of line numbers
numberstyle=\tiny\color{purple!80!white},   % Style of line numbers
numbersep=1em,
backgroundcolor=\color{gray!10}, % Background color
xleftmargin=0.05\linewidth,                % Left margin
xrightmargin=-2pt,
aboveskip=10pt,
belowskip=10pt,
frame=l,                    % Frame around the code
rulecolor=\color{black!40},
framerule=0.1pt,
breaklines=false,                 % Allow line breaks
showstringspaces=false,          % Don't show spaces in strings
linewidth=0.95\linewidth,            % Fixed width size
basicstyle=\ttfamily\footnotesize, % Font style and size
captionpos=t,                     % Caption position
language=bash,
literate=
{┐}{\textSFiii{}}1%
{└}{\textSFii{}}1%
{┴}{\textSFvii{}}1%
{┬}{\textSFvi{}}1%
{├}{\textSFviii{}}{1}%
{─}{\textSFx{}}1%
{┼}{\textSFv{}}1%
{│}{\textSFxi{}}1%
{~}{{\mtildeascii}}1
{🙂}{\emoji {🙂}}1%      
}

\newtcblisting{foo}{
listing only,
nobeforeafter,
after={\xspace},
hbox,
tcbox raise base,
fontupper=\ttfamily,
colback=lightgray,
colframe=lightgray,
size=fbox
}

\definecolor{lightgray}{RGB}{240, 240, 240}
\definecolor{highlightcolor}{RGB}{255, 142, 143}
\definecolor{correctcolor}{RGB}{167, 201, 87}

\newcommand{\hb}[1]{\scalebox{0.75}{\colorbox{highlightcolor}{\scalebox{1.2}{#1}}}}
\newcommand{\hcb}[1]{\scalebox{0.75}{\colorbox{correctcolor}{\scalebox{1.2}{#1}}}}


\usepackage{newunicodechar}

%% \DeclareUnicodeCharacter{2320}{/}
%% \DeclareUnicodeCharacter{23AE}{|}
%% \DeclareUnicodeCharacter{2321}{/}
%% \DeclareUnicodeCharacter{2572}{\textbackslash}
%% \DeclareUnicodeCharacter{2571}{/}
%% \DeclareUnicodeCharacter{3C6}{$\phi$}
%% \DeclareUnicodeCharacter{2080}{${}_0$}
%% \DeclareUnicodeCharacter{3C0}{$\pi$}
%% \DeclareUnicodeCharacter{3A9}{$\Omega$}
%% %% \DeclareUnicodeCharacter{3A9}{🙂}
%% \newunicodechar{🙂}{\texttt{<smile>}}

%% Closes the \makeatletter opened before the listing-badge block above; from
%% here on `@' is an ordinary character again.
\makeatother

%% ---------------------------------------------------------------------------
%% Inline code chips.
%%
%% \token typesets a fragment of Ymir inside a tinted box, in the middle of a
%% sentence. Two things had to be fixed for it to sit in running text:
%%
%%  - \fboxsep defaults to 3pt, so every chip carried 3pt of padding on each
%%    side. Against a following comma or full stop that padding read as a
%%    space -- "the function println , which" -- because the punctuation was
%%    pushed a visible gap away from the word it belongs to. The padding is now
%%    1.2pt, enough to keep the tint off the glyphs and no more.
%%  - the chip was set at \footnotesize, two steps below the surrounding 11pt
%%    body, which made inline identifiers noticeably smaller than the prose
%%    naming them. It is now \small.
%%
%% A \colorbox is an unbreakable box, so a long chip cannot be hyphenated and
%% has only the two line positions the surrounding words leave it. What buys
%% the room to place it is \emergencystretch (set with the vertical rhythm
%% further down), not an \allowbreak here: a break permitted at the very start
%% of a chip opens an empty first line whenever the chip begins a narrow table
%% cell.
\definecolor{tokenbg}{rgb}{0.945, 0.930, 0.902}

\newcommand{\token}[1]{%
  \begingroup
  \setlength{\fboxsep}{1.2pt}%
  \colorbox{tokenbg}{\lstinline[style=coloredverbatim,%
    basicstyle=\ttfamily\small]{#1}}%
  \endgroup}

\newcommand{\tokennolst}[1]{%
  \begingroup
  \setlength{\fboxsep}{1.2pt}%
  \colorbox{tokenbg}{\ttfamily\small#1}%
  \endgroup}


\usepackage{graphicx}
\usepackage{longtable}
\usepackage{wrapfig}
\usepackage{rotating}
\usepackage[normalem]{ulem}
\usepackage{amsmath}
\usepackage{amssymb}
\usepackage{capt-of}
%% Links. The default is a red rectangle drawn around every link, which in a
%% book whose table of contents is entirely links means a page of red boxes.
%% Internal references are instead tinted a dark, desaturated blue: legible as
%% a link on screen, and close enough to black to disappear in a print.
\definecolor{linkcolor}{rgb}{0.13, 0.20, 0.42}
\definecolor{urlcolor}{rgb}{0.10, 0.35, 0.42}
\usepackage[colorlinks=true,
            linkcolor=linkcolor,
            citecolor=linkcolor,
            urlcolor=urlcolor,
            linktoc=all]{hyperref}

\tikzstyle{every picture}+=[remember picture]
\tikzstyle{mProjectPP}=[circle,fill=yellow!70, minimum size=20pt,inner sep=0pt]
\tikzstyle{label}=[rectangle, fill=white, minimum size=40pt, inner sep=0pt]
\tikzstyle{Box}=[rectangle, fill=white, minimum size=40pt, inner sep=0pt]
\tikzstyle{Box2}=[rectangle, fill=white, minimum size=20pt, inner sep=0pt]


\usepackage{caption}
\usepackage{dblfloatfix}


\usepackage{makecell}

\setlength{\columnsep}{1cm}
\setlength{\columnseprule}{0.05cm}
\def\columnseprulecolor{\color{black!10}}

\usetikzlibrary {arrows.meta, fit, calc, arrows, positioning}

%% A listing caption is a rule the width of the band of colour below it, with
%% the caption text sitting on that rule -- so the caption, the rule and the
%% code read as one object. Both are indented by the listings' own
%% `xleftmargin' (0.05\linewidth); they used to start at 0.02 and end at 0.93,
%% aligning with neither the code band nor the text block.
\DeclareCaptionFormat{listing}{%
  \hspace*{0.05\linewidth}{\color{black!45}\rule{0.95\linewidth}{0.8pt}}\par
  \hspace*{0.05\linewidth}#1#2#3\vspace{1pt}}
\captionsetup[lstlisting]{format=listing, singlelinecheck=false,
  margin=0pt, font={sf,small}, labelfont={sf,small,bf}}

%% Figure and table captions join the listing captions in the sans face, so
%% that everything labelling a floating object looks alike.
\captionsetup[figure]{font={sf,small}, labelfont={sf,small,bf}}
\captionsetup[table]{font={sf,small}, labelfont={sf,small,bf}}
%\renewcommand\lstlistingname{Example}

\usepackage{awesomebox}
\usepackage{subcaption}

\usetikzlibrary{chains,shadows.blur}
\usepackage{afterpage}
\usepackage{pagecolor}

\usepackage{siunitx}

%% ---------------------------------------------------------------------------
%% Callout boxes.
%%
%% The book used to have two visual treatments for the same thing: awesomebox's
%% own \notebox, \warningbox and \importantbox, which put an icon and a rule
%% beside plain text, and locally defined \mynotebox, \mytipbox and
%% \mycautionbox, which additionally laid a grey tint behind it. Which one an
%% aside got was a matter of which chapter it was written in. All five are now
%% tinted, so an aside looks like an aside wherever it appears, and the icon
%% alone says what kind it is.
%%
%% The minipage has to leave room for the padding \colorbox puts around it, or
%% the tint -- and the text inside it -- overruns the box by 2\fboxsep.
\newcommand{\ymirtintedbox}[1]{%
  \colorbox{gray!10!white}{%
    \begin{minipage}{\dimexpr\linewidth-2\fboxsep\relax}#1\end{minipage}}}

\makeatletter
%% Wrap an existing callout command so that its content is tinted, keeping the
%% original under a private name.
\newcommand{\ymir@tintbox}[1]{%
  \expandafter\let\csname ymir@plain#1\expandafter\endcsname\csname #1\endcsname
  \expandafter\renewcommand\csname #1\endcsname[1]{%
    \csname ymir@plain#1\endcsname{\ymirtintedbox{##1}}}}
\ymir@tintbox{notebox}
\ymir@tintbox{tipbox}
\ymir@tintbox{cautionbox}
\ymir@tintbox{warningbox}
\ymir@tintbox{importantbox}
\makeatother

\newcommand{\fileicon}[4][gray!10]{%
% #1 = fill color (optional, default gray!10)
% #2 = x position
% #3 = y position
% #4 = label text
\begin{scope}[shift={({#2},{#3})}]
\def\filewidth{2cm}
\def\fileheight{2.5cm}
\def\foldsize{0.5cm}

% Base file shape (5 sides)
\filldraw[fill=#1, draw=black, thick]
(0,0) --
(0,\fileheight) --
(\filewidth-\foldsize,\fileheight) --
(\filewidth,\fileheight-\foldsize) --
(\filewidth,0) -- cycle;

% Folded flap
\filldraw[fill=gray!30, draw=black, thick]
(\filewidth-\foldsize,\fileheight) --
(\filewidth,\fileheight-\foldsize) --
(\filewidth-\foldsize,\fileheight-\foldsize) -- cycle;

% Label
\node[font=\sffamily\bfseries, align=center, text=black]
at (\filewidth/2, \fileheight/2 - 0.2cm) {#4};
\end{scope}
}


\newcommand{\foldericon}[4][yellow!60]{%
% #1 = fill color (default yellow!60)
% #2 = x position
% #3 = y position
% #4 = label
\begin{scope}[shift={({#2},{#3})}]
\def\fw{3cm}   % folder width
\def\fh{1.8cm} % folder height
\def\tw{0.4cm} % tab width
\def\th{0.2cm} % tab height
\def\radius{0.1cm}

% Folder body (rounded rectangle, flat)
\filldraw[fill=#1, draw=black, line width=0.5pt, rounded corners=\radius]
(0,0) -- (\fw,0) -- (\fw,\fh-\th) -- (0,\fh-\th) -- cycle;

% Tab (flap)
\filldraw[fill=#1, draw=black, line width=0.5pt, rounded corners=\radius]
(0,\fh-\th) -- (\tw,\fh) -- (\tw+0.6cm,\fh) -- (\tw+0.6cm,\fh-\th) -- cycle;

\node[font=\sffamily\bfseries, text=black]
at (\fw/2, \fh/2 - 0.2cm) {#4};

%% % Label
%% \node[font=\sffamily\bfseries, text=black]
%%   at (\fw/2, (\fh-\th)/2) {#4};

\end{scope}
}


% --- Gear icon with trapezoid teeth (flat-top) ---
\newcommand{\gearicon}[9][black]{%
% #1 = color
% #2 = x position
% #3 = y position
% #4 = number of teeth
% #5 = gear radius
% #6 = tooth height
% #7 = tooth base half-angle (deg)
% #8 = tooth top half-angle (deg)

\begin{scope}[shift={({#2},{#3})}, rotate=#4]
\def\radius{#5}
\def\toothheight{#6}
\def\baseangle{#7}
\def\topangle{#8}
\def\nteeth{7}

% Draw each tooth
\foreach \i in {0,...,\numexpr\nteeth-1\relax} {
\begin{scope}[rotate={360 /\nteeth*\i}]
\draw[fill=#1, draw=#1]
({\radius*cos(\baseangle)}, {\radius*sin(\baseangle)}) --        % left base
({(\radius+\toothheight)*cos(\topangle)}, {(\radius+\toothheight)*sin(\topangle)}) --  % left top
({(\radius+\toothheight)*cos(-\topangle)}, {(\radius+\toothheight)*sin(-\topangle)}) -- % right top
({\radius*cos(-\baseangle)}, {\radius*sin(-\baseangle)}) --       % right base
cycle;
\end{scope}
}

% Gear body
\fill[fill=#1] (0,0) circle (\radius);

% Center hole
\fill[fill=white] (0,0) circle (\radius/#9);
\end{scope}
}




% --- Command for rotated rounded rectangle ---
\newcommand{\chainicon}[8][black]{%
% #1 = color
% #2 = x position
% #3 = y position
% #4 = width
% #5 = height
% #6 = line width

%% \scalebox{#7}{
\begin{scope}[shift={({#2},{#3})}, rotate=#8]

\draw[line width=#6, color=#1!70, rounded corners=0.15cm, fill=none]
(-#4/2,-#5/2) rectangle (#4/2,#5/2);

\draw[line width=#6/2.8, color=#1!80, rounded corners=0.15cm, fill=none]
(-#4/2,-#5/2) rectangle (#4/2,#5/2);

\draw[line width=#6, color=#1!70, rounded corners=0.15cm, fill=none]
(-#4/2+#4*1.25,-#5/2) rectangle (#4/2+#4*1.25,#5/2);

\draw[line width=#6/2.8, color=#1!80, rounded corners=0.15cm, fill=none]
(-#4/2+#4*1.25,-#5/2) rectangle (#4/2+#4*1.25,#5/2);

\draw[line width=#6+0.03cm, color=white, rounded corners=0.05cm]
(-#4/2+#4*0.625,-0.01) rectangle (#4/2+#4*0.625,0.01);

\draw[line width=#6, color=#1!70, rounded corners=0.0025cm]
(-#4/2+#4*0.625,-0.01) rectangle (#4/2+#4*0.625,0.01);

\draw[line width=#6/2.8, color=#1!80, rounded corners=0.0025cm]
(-#4/2+#4*0.625,-0.01) rectangle (#4/2+#4*0.625,0.01);

\end{scope}
%}
}

\newcommand{\hammer}[5][black] {
% #1 = color
% #2 = x position
% #3 = y position
% #4 = width

\begin{scope}[shift={({#2},{#3})}]
\scalebox{#5}{
\draw[line width=0.1, color=brown!30, rounded corners=0.02cm, fill=brown!70, rotate=#4]
plot coordinates {(0, 0) (0.025, 0.5) (0, 0.7) (0.2,0.7) (0.1725, 0.5) (0.2,0)} -- cycle;

\draw[line width=0.001cm, color=gray!30, rounded corners=0.05cm, fill=gray!70, rotate=#4]
(-0.2, 0.6) rectangle (0.4, 0.9);

%% \scalebox{0.8}{
\node[font=\sffamily\bfseries, text=black]
at (0.85, 0.4) {\fontsize{1pt}{1pt}\selectfont \{\}};
%%}

}

\end{scope}
}

%% ---------------------------------------------------------------------------
%% Vertical rhythm.
%%
%% The book sets paragraphs block-style: no first-line indent, blank space
%% between paragraphs instead. That space used to be a rigid 10pt, which gave
%% TeX nothing to stretch or shrink when filling a page, so it had to resort to
%% breaking pages badly (and to the stray single line stranded at the top or
%% bottom of a page). It is now elastic, and widows and orphans are forbidden
%% outright.
\setlength{\parindent}{0pt}
\setlength{\parskip}{0.6\baselineskip plus 3pt minus 2pt}

%% Termes, like the Times it clones, is a narrow face with tight default
%% leading; opening it slightly is what makes a long paragraph readable.
\linespread{1.06}

\widowpenalty=10000
\clubpenalty=10000
\displaywidowpenalty=10000
%% ... but allow a page to run a line short rather than stretch it to death.
\raggedbottom

%% Character protrusion and font expansion. Nothing else buys as much evenness
%% of colour in a justified text block for as little effort.
\usepackage{microtype}

%% \token typesets an inline code fragment inside a \colorbox, and a box cannot
%% be hyphenated. A long one -- \token{foo::bar::pubFuncInBar} and its like --
%% therefore has exactly two possible positions in a line, and when neither fits
%% TeX gives up and lets it run into the margin. \emergencystretch grants a
%% final pass with extra stretch between words, which buys the room to push the
%% chip onto the next line instead.
\setlength{\emergencystretch}{3em}
\tolerance=1200

%% Tables. The wider ones declare their columns as fixed fractions of
%% \textwidth summing to nearly 1, which leaves nothing for the padding LaTeX
%% adds on each side of every column: at the default 6pt they all ran a couple
%% of millimetres past the text block. 3pt still separates the columns clearly.
%% The extra leading is what booktabs rules want around them.
\setlength{\tabcolsep}{3pt}
\renewcommand{\arraystretch}{1.15}

%% ---------------------------------------------------------------------------
%% Part, chapter and section heads.
%%
%% The body is Termes (a Times clone) and the furniture -- figure labels,
%% listing badges, captions -- is Heros (a Helvetica clone). The heads join the
%% furniture rather than the body: sans-serif heads over a seriffed text block
%% is the standard technical-book pairing, and it makes a head unmistakably a
%% head at a glance rather than just bigger prose.
\usepackage{titlesec}

%% A hairline closing a head block.
%%
%% Built from a bare \hrule rather than titlesec's \titlerule: \titlerule sets
%% the rule in horizontal mode, which left it sitting on the baseline of the
%% first line of the following paragraph instead of on its own line. \hrule is
%% a vertical-mode primitive and contributes its own height to the page.
\newcommand{\ymirheadrule}[1][0.8pt]{%
  \par\nobreak\vskip 5pt%
  {\color{black!25}\hrule height #1}%
  \par\nobreak\vskip 12pt}

%% Part: the number spelled out small above a large title, over a full rule.
\titleformat{\part}[display]
  {\normalfont\sffamily\filcenter}
  {\fontsize{13}{15}\selectfont\color{black!45}%
   \MakeUppercase{\partname}\ \thepart}
  {14pt}
  {\fontsize{30}{34}\selectfont\bfseries}
  %% Braced: the optional argument of \ymirheadrule would otherwise end
  %% titlesec's own optional argument at its first `]'.
  [{\ymirheadrule[1.2pt]}]
\titlespacing*{\part}{0pt}{0pt}{28pt}

%% Chapter: "CHAPTER n" in small caps-height sans, the number set large beside
%% it, the title flush right beneath, and a hairline closing the block.
\titleformat{\chapter}[display]
  {\normalfont\sffamily\raggedleft}
  {\fontsize{12}{14}\selectfont\color{black!45}%
   \MakeUppercase{\chaptertitlename}\hspace{0.35em}%
   {\fontsize{40}{40}\selectfont\color{black!70}\thechapter}}
  {10pt}
  {\fontsize{24}{28}\selectfont\bfseries}
  [\ymirheadrule]
\titlespacing*{\chapter}{0pt}{20pt}{34pt}

%% Unnumbered chapters (Preamble, Contents) get the same block without the
%% number, so they do not read as a different kind of division.
\titleformat{name=\chapter,numberless}[display]
  {\normalfont\sffamily\raggedleft}
  {}
  {0pt}
  {\fontsize{24}{28}\selectfont\bfseries}
  [\ymirheadrule]
\titlespacing*{name=\chapter,numberless}{0pt}{20pt}{34pt}

\titleformat{\section}
  {\normalfont\sffamily\fontsize{15}{18}\selectfont\bfseries}
  {\color{black!55}\thesection}{0.7em}{}
\titlespacing*{\section}{0pt}{1.6\baselineskip plus 4pt minus 2pt}{0.5\baselineskip}

\titleformat{\subsection}
  {\normalfont\sffamily\fontsize{12.5}{15}\selectfont\bfseries}
  {\color{black!55}\thesubsection}{0.7em}{}
\titlespacing*{\subsection}{0pt}{1.2\baselineskip plus 3pt minus 2pt}{0.35\baselineskip}

\titleformat{\subsubsection}
  {\normalfont\sffamily\fontsize{11}{13}\selectfont\bfseries}
  {\color{black!55}\thesubsubsection}{0.6em}{}
\titlespacing*{\subsubsection}{0pt}{1\baselineskip plus 2pt minus 2pt}{0.3\baselineskip}

%% ---------------------------------------------------------------------------
%% Running heads: page number to the outer edge, the division name to the
%% inner, both small sans over a hairline. The class default sets the chapter
%% title in italic capitals across the whole head, which at this size competes
%% with the text below it.
\usepackage{fancyhdr}
\pagestyle{fancy}
\fancyhf{}
\renewcommand{\headrulewidth}{0.4pt}
\renewcommand{\footrulewidth}{0pt}
\renewcommand{\headrule}{%
  {\color{black!30}\hrule height\headrulewidth width\headwidth}}
\fancyhead[LE]{\sffamily\footnotesize\color{black!70}\thepage}
\fancyhead[RE]{\sffamily\footnotesize\color{black!55}\nouppercase{\leftmark}}
\fancyhead[LO]{\sffamily\footnotesize\color{black!55}\nouppercase{\rightmark}}
\fancyhead[RO]{\sffamily\footnotesize\color{black!70}\thepage}

%% A chapter's own opening page carries no running head, only its number.
\fancypagestyle{plain}{%
  \fancyhf{}%
  \renewcommand{\headrulewidth}{0pt}%
  \fancyfoot[C]{\sffamily\footnotesize\color{black!70}\thepage}}

\renewcommand{\chaptermark}[1]{\markboth{\thechapter.\ #1}{}}
\renewcommand{\sectionmark}[1]{\markright{\thesection.\ #1}}

\usepackage{minitoc}
\setcounter{minitocdepth}{5}

%% The per-chapter mini table of contents is a navigation aid, not a place for
%% a page of coloured link text: its entries stay black.
\mtcsetfeature{minitoc}{before}{\begingroup\hypersetup{linkcolor=black}}
\mtcsetfeature{minitoc}{after}{\endgroup}

%% The mini table of contents is furniture: small sans, flush with the text
%% block rather than indented a quarter of the way across it.
\setlength{\mtcindent}{0pt}
\mtcsetfont{minitoc}{section}{\sffamily\small}
\mtcsetfont{minitoc}{subsection}{\sffamily\small\color{black!75}}
\mtcsetfont{minitoc}{subsubsection}{\sffamily\small\color{black!75}}
\mtcsettitlefont{minitoc}{\sffamily\normalsize\bfseries}

\usepackage{tocloft}
\setlength{\cftparskip}{0pt}

\renewcommand{\mtcskip}{0pt}  % spacing between entries

%% A table of contents is a list of heads, so it is set in the same sans face
%% as the heads themselves rather than in the body face.
\renewcommand{\cfttoctitlefont}{\hfill\normalfont\sffamily\fontsize{24}{28}\selectfont\bfseries}
\renewcommand{\cftaftertoctitle}{\ymirheadrule}
\renewcommand{\cftpartfont}{\sffamily\large\bfseries}
\renewcommand{\cftchapfont}{\sffamily\bfseries}
\renewcommand{\cftsecfont}{\sffamily}
\renewcommand{\cftsubsecfont}{\sffamily\small\color{black!75}}
\renewcommand{\cftchappagefont}{\sffamily\bfseries}
\renewcommand{\cftsecpagefont}{\sffamily}
\renewcommand{\cftsubsecpagefont}{\sffamily\small\color{black!75}}

%% Leader dots at every level of the main table of contents were dense enough
%% to read as a grey rule; they are thinned out and lightened here.
\renewcommand{\cftdotsep}{4}
\renewcommand{\cftchapleader}{\color{black!35}\cftdotfill{\cftdotsep}}
\renewcommand{\cftsecleader}{\color{black!35}\cftdotfill{\cftdotsep}}
\renewcommand{\cftsubsecleader}{\color{black!35}\cftdotfill{\cftdotsep}}

%% A page left blank to start a chapter on a recto should be blank: the class
%% default still prints the running head and the page number on it.
\renewcommand{\cleardoublepage}{%
  \clearpage
  \ifodd\value{page}\else
    \thispagestyle{empty}\hbox{}\newpage
  \fi}

%% ---------------------------------------------------------------------------
%% Title page.
%%
%% \maketitle gave the class default: the title, the author and the date
%% centred in the middle of an otherwise empty page, in the body face, with no
%% indication of what the book contains. This one is set in the same sans face
%% as the heads, anchored to a rule near the top of the page, and says what the
%% two parts of the book are.
\makeatletter
\newcommand{\ymirtitlepage}{%
  \begin{titlepage}
    \thispagestyle{empty}
    \sffamily
    \raggedright
    \vspace*{0.18\textheight}

    {\color{black!30}\rule{\linewidth}{1.2pt}}

    \vspace{1.2em}
    {\fontsize{15}{18}\selectfont\color{black!55}
      THE YMIR PROGRAMMING LANGUAGE\par}

    \vspace{0.9em}
    {\fontsize{34}{40}\selectfont\bfseries\@title\par}

    \vspace{1.4em}
    {\color{black!30}\rule{\linewidth}{0.5pt}}

    \vspace{1.2em}

  \end{titlepage}}
\makeatother

\newcommand{\forceevenpage}{%
  \ifodd\value{page}\hbox{}\newpage\fi
}
